%----------------------------------------------------------------------
\section{Experiments}
\label{sec:experiments}
%----------------------------------------------------------------------

The experiments answer four questions in order: what accuracy the framework reaches under a small
time-series pretraining budget (\S\ref{sec:main}, \S\ref{sec:budget}); what the video initialisation
contributes and where that contribution stops (\S\ref{sec:backbone}); what the transfer design
contributes (\S\ref{sec:design}); and why the naive alternative fails (\S\ref{sec:diagnosis}).

\subsection{Setup}
\label{sec:setup}

\paragraph{Protocol.} We follow the zero-shot long-sequence protocol of \citet{visionts}:
Time-Series-Library splits \citep{tslib} (ETT 12/4/4 months, others 70/10/20), a
\texttt{StandardScaler} fitted on the target dataset's training portion, all channels, stride-1 test
origins, $H\in\{96,192,336,720\}$, MSE/MAE on standardised values, and their per-dataset context
lengths and periods (Appendix~\ref{app:impl}). Running their released checkpoint through our harness
reproduces their published numbers to three decimals (ETTh1 at $H{=}96$: $0.3527/0.3833$ against
$0.353/0.383$). Zero-shot here means no parameter is updated on any evaluation dataset and no
evaluation dataset appears in the pretraining corpus; the model is trained once and applied
unchanged, with one checkpoint and one inference rule for all 24 cells.

\paragraph{Ablations and baselines.} Ablations and initialisation studies use five datasets
(ETT$\times$4 and weather) at three horizons, which is what one evaluation budget covers; the
six-dataset four-horizon grid is reserved for the main table. Five- and six-dataset means are never
mixed and every table says which it reports. Baseline numbers are quoted from \citet{visiontspp}
and \citet{visionts}; we did not retrain them.

\subsection{Accuracy under a small continual-pretraining budget}
\label{sec:main}

% generated by paper/make_main_table.py -- do not edit by hand
\begin{table}[t]
\caption{Zero-shot long-sequence forecasting. Each entry averages the four horizons $H\in\{96,192,336,720\}$; MSE and MAE are computed on standardised values. Our model uses one checkpoint and one inference rule for all 24 cells. Baselines are quoted from \citet{visiontspp} and \citet{visionts}; `--' means the source does not report that dataset. \textbf{Bold}: best. \underline{Underlined}: second best. The last column counts first places over the twelve dataset columns.}
\label{tab:main}
\begin{center}
\setlength{\tabcolsep}{3pt}
\scriptsize
\begin{tabular}{lccccccccccccccc}
\toprule
\multirow{2}{*}{Model} & \multicolumn{2}{c}{ETTm1} & \multicolumn{2}{c}{ETTm2} & \multicolumn{2}{c}{ETTh1} & \multicolumn{2}{c}{ETTh2} & \multicolumn{2}{c}{ECL} & \multicolumn{2}{c}{Weather} & \multicolumn{2}{c}{\textbf{Avg}} & \multirow{2}{*}{1\textsuperscript{st}} \\
\cmidrule(lr){2-3} \cmidrule(lr){4-5} \cmidrule(lr){6-7} \cmidrule(lr){8-9} \cmidrule(lr){10-11} \cmidrule(lr){12-13} \cmidrule(lr){14-15}
 & MSE & MAE & MSE & MAE & MSE & MAE & MSE & MAE & MSE & MAE & MSE & MAE & MSE & MAE & \\
\midrule
\textbf{Ours} & \textbf{0.352} & \textbf{0.369} & \textbf{0.241} & \underline{0.302} & 0.404 & \underline{0.416} & 0.335 & 0.376 & \textbf{0.163} & \textbf{0.254} & \textbf{0.215} & 0.252 & \textbf{0.285} & 0.328 & 6 \\
VisionTS++ (ViT-L) & \underline{0.354} & \textbf{0.369} & \underline{0.244} & \textbf{0.298} & 0.403 & 0.418 & \textbf{0.327} & \textbf{0.365} & \underline{0.181} & \underline{0.264} & 0.226 & \underline{0.243} & \underline{0.289} & \textbf{0.326} & 4 \\
VisionTS++ (ViT-B) & 0.360 & \underline{0.372} & \underline{0.244} & \textbf{0.298} & 0.402 & \underline{0.416} & \underline{0.333} & \underline{0.370} & 0.184 & 0.265 & \underline{0.222} & \textbf{0.241} & 0.291 & \underline{0.327} & 2 \\
VisionTS & 0.374 & \underline{0.372} & 0.318 & 0.366 & \textbf{0.390} & \textbf{0.414} & \underline{0.333} & 0.375 & 0.207 & 0.294 & 0.269 & 0.292 & 0.315 & 0.352 & 2 \\
\midrule
Moirai$_\mathrm{S}$ & 0.448 & 0.410 & 0.272 & 0.321 & 0.400 & 0.424 & 0.341 & 0.379 & 0.233 & 0.320 & 0.242 & 0.267 & 0.323 & 0.353 & -- \\
Moirai$_\mathrm{B}$ & 0.382 & 0.388 & 0.276 & 0.320 & 0.434 & 0.439 & 0.346 & 0.382 & 0.188 & 0.274 & 0.238 & 0.261 & 0.311 & 0.344 & -- \\
Moirai$_\mathrm{L}$ & 0.390 & 0.389 & 0.317 & 0.366 & 0.510 & 0.469 & 0.354 & 0.377 & 0.188 & 0.273 & 0.260 & 0.275 & 0.337 & 0.358 & -- \\
Chronos$_\mathrm{S}$ & 0.640 & 0.500 & 0.310 & 0.350 & 0.545 & 0.472 & 0.424 & 0.430 & 0.220 & 0.284 & 0.300 & 0.318 & 0.407 & 0.392 & -- \\
Chronos$_\mathrm{B}$ & 0.646 & 0.500 & 0.295 & 0.338 & 0.591 & 0.468 & 0.406 & 0.411 & 0.215 & 0.279 & 0.293 & 0.315 & 0.408 & 0.385 & -- \\
Chronos$_\mathrm{L}$ & 0.556 & 0.465 & 0.300 & 0.341 & 0.589 & 0.466 & 0.455 & 0.427 & 0.204 & 0.274 & 0.279 & 0.306 & 0.397 & 0.380 & -- \\
Time-MoE$_\mathrm{S}$ & 0.394 & 0.416 & 0.316 & 0.361 & 0.400 & 0.424 & 0.367 & 0.404 & -- & -- & 0.266 & 0.297 & -- & -- & -- \\
Time-MoE$_\mathrm{B}$ & 0.376 & 0.406 & 0.349 & 0.380 & \underline{0.394} & 0.420 & 0.405 & 0.415 & -- & -- & 0.270 & 0.300 & -- & -- & -- \\
Timer & 0.487 & 0.457 & 0.328 & 0.347 & 0.444 & 0.457 & 0.358 & 0.407 & -- & -- & 0.304 & 0.331 & -- & -- & -- \\
TimesFM & 0.433 & 0.419 & -- & -- & 0.473 & 0.444 & 0.392 & 0.406 & -- & -- & -- & -- & -- & -- & -- \\
MOMENT & 0.670 & 0.537 & 0.316 & 0.371 & 0.684 & 0.566 & 0.362 & 0.410 & 0.765 & 0.687 & 0.294 & 0.326 & 0.515 & 0.483 & -- \\
\bottomrule
\end{tabular}
\end{center}
\end{table}


Table~\ref{tab:main} reports four-horizon averages. Our model attains the lowest mean MSE among the
models compared on this protocol, $0.285$, ahead of VisionTS++ large ($0.289$) and base ($0.291$),
and takes the most first places over the twelve dataset columns. The gains are concentrated in
electricity ($0.163$ against $0.181$) and weather ($0.215$ against $0.222$). Two datasets go the
other way: ETTh1 to VisionTS ($0.390$ against our $0.404$) and ETTh2 to VisionTS++ large ($0.327$
against $0.335$), both hourly sets where a longer tuned context helps at long horizons. Mean MAE is
$0.328$ against VisionTS++ large's $0.326$: we are behind on that metric, and we have no uncertainty
estimate that would let us call the difference immaterial.

At 20k updates instead of 60k the same recipe reaches $0.288/0.331$, still first on mean MSE here at
a third of the training.


Averaged over the six datasets, error grows from $0.223/0.281$ at $H{=}96$ to $0.349/0.377$ at
$H{=}720$; the full per-horizon breakdown is Table~\ref{tab:perhorizon} in
Appendix~\ref{app:perhorizon}, and the two datasets we lose overall are lost at the mid-to-long
horizons rather than uniformly.

\subsection{Pretraining budget}
\label{sec:budget}

% Budget comparison. Every column is a quantity each paper actually reports; nothing is converted
% into FLOPs or wall-clock across systems.
\begin{table}[t]
\caption{Time-series pretraining budget in the units each paper reports. Sample presentations are
optimizer updates $\times$ effective batch, i.e.\ how many rendered windows the model was shown;
they are not a compute measure, since the two systems differ in model size and in the number of time
points and visual tokens per sample. The upstream video or image pretraining that both inherit is
excluded, as is any FLOP or wall-clock comparison across systems. VisionTS++ figures are as reported
in \citet{visiontspp}. Takeaway: the reported accuracy is obtained from $\approx$0.3\% of the
archive and $\approx$1/27 of the sample presentations, in 18.9 hours on one GPU.}
\label{tab:budget}
\begin{center}
\setlength{\tabcolsep}{5pt}
\footnotesize
\begin{tabular}{lccc}
\toprule
 & Ours (main) & Ours (20k) & VisionTS++ \\
\midrule
Time-series corpus (observations) & 654.2M & 654.2M & LOTSA archive, 231B \\
\quad reachable after the 5\% holdout & 621.9M & 621.9M & not reported \\
Optimizer updates & 60{,}000 & 20{,}000 & 100{,}000 \\
Effective batch (windows/update) & 32 & 32 & 512 \\
Sample presentations & 1.92M & 0.64M & 51.2M \\
Backbone & ViT-B (87M) & ViT-B (87M) & ViT-B / ViT-L \\
Effective passes over the corpus & $\approx$6.8 & $\approx$2.3 & not reported \\
Measured training time & 18.9\,h, 1$\times$A40 & 6.3\,h, 1$\times$A40 & not reported \\
\bottomrule
\end{tabular}
\end{center}
\end{table}


Table~\ref{tab:budget} states the resources on both sides in the units each paper reports. We use a
corpus of 654.2M observations and 60k optimizer updates at an effective batch of 32. We report
corpus size, optimizer updates and sample presentations separately, and we do not convert them into
a compute claim: the models differ in size, the number of time points and visual tokens per sample
differs between the two renderings, and we do not know the baselines' hardware or wall-clock. The
comparison also excludes the original video pretraining of our backbone, which we inherit rather
than pay for, exactly as the image-based systems inherit ImageNet pretraining.

\subsection{What the video initialisation contributes}
\label{sec:backbone}

% Initialisation study at two adaptation budgets: the 5k advantage and the 20k catch-up together.
\begin{table}[t]
\caption{Initialisation study. Rendering, objective, corpus, schedule and inference rule are held
fixed; only the weights the model starts from change. Five-dataset mean MSE (ETT$\times$4 and
weather) at $H{=}96/336/720$. Two seeds are listed where both were run. At a 5k-update budget both
video initialisations lead the image initialisation; by 20k the image initialisation has caught up,
which is why we claim adaptation efficiency rather than a higher ceiling. The V-JEPA arms are
size-matched, not architecturally identical (see text).}
\label{tab:backbone}
\begin{center}
\setlength{\tabcolsep}{4pt}
\footnotesize
\begin{tabular}{llcll}
\toprule
Initialisation & Pretraining signal & Params & 5k updates & 20k updates \\
\midrule
ImageNet MAE-B & images & 87M & 0.264 / 0.350 / 0.395
  & 0.244 / 0.329 / \textbf{0.373} \\
 & & & & 0.245 / 0.330 / \textbf{0.373} \\
VideoMAE-B & video, reconstructive & 87M & \textbf{0.250} / 0.337 / 0.384
  & 0.245 / 0.330 / 0.379 \\
V-JEPA~2.1-B & video, predictive & 87M\,+\,23M & 0.252 / \textbf{0.335} / \textbf{0.382}
  & not run \\
\midrule
V-JEPA~2-L & video, predictive & 300M & 0.245 / 0.329 / 0.379 & not run \\
V-JEPA~2-L, random init & none & 300M & 0.280 / 0.369 / 0.418 & not run \\
\bottomrule
\end{tabular}
\end{center}
\end{table}


Table~\ref{tab:backbone} holds the rendering, objective, corpus, schedule and inference rule fixed
and changes only the weights the model starts from. At a 5k-step adaptation budget, both video
initialisations lead the ImageNet initialisation by $3$--$5\%$ at every horizon, and they are within
$0.002$ of each other, so the reconstructive and predictive video objectives are not distinguishable
here. The random-init control is the informative row: the same ViT-L architecture, same data, same
schedule, reaches $0.280$ at $H{=}96$ against $0.245$ for its pretrained twin, and is worse than an
image-pretrained backbone $3.5\times$ smaller. Whatever transfers is not capacity.

Two qualifications belong with this table. First, the V-JEPA arms are size-matched rather than
architecturally identical: they carry an additional predictor, use EMA targets, and their predictor
output head is re-initialised because the pretrained head maps into a teacher space our targets do
not use. Second, the advantage is budget-dependent. The right-hand column runs the same comparison
at 20k updates, where the image initialisation has caught up and is level or slightly ahead at the
longest horizon, so what the video prior buys is adaptation efficiency and not a higher ceiling.
Note also that the image-initialised arm is an initialisation control, not a reimplementation of
VisionTS++, which is a different system with a different rendering and a multi-quantile training
objective.

\subsection{What the transfer design contributes}
\label{sec:design}

% Objective ablation and budget curve under the forecast-aligned objective.
\begin{table}[t]
\caption{Ablation of the transfer design. Five-dataset mean MSE (ETT$\times$4 and weather),
32-frame clips, identical corpus, masking, schedule and inference rule; only the listed component
changes.}
\label{tab:ablation}
\begin{center}
\setlength{\tabcolsep}{6pt}
\footnotesize
\begin{tabular}{lccc}
\toprule
Configuration & $H{=}96$ & $H{=}336$ & $H{=}720$ \\
\midrule
Full framework, 20k updates & 0.245 & 0.330 & 0.379 \\
\quad pixel-only supervision (no value-space term) & 0.257 & 0.347 & 0.411 \\
\quad no prior anchoring & 0.244 & 0.330 & 0.378 \\
\midrule
Full framework, 5k updates & 0.250 & 0.337 & 0.384 \\
Full framework, 60k updates & \textbf{0.242} & \textbf{0.329} & \textbf{0.377} \\
\bottomrule
\end{tabular}
\end{center}
\end{table}


Table~\ref{tab:ablation} varies the training objective and the budget with everything else fixed.
Replacing the forecast-aligned value-space term with pixel-only supervision, holding masking,
rendering, corpus and schedule constant, costs $5\%$ at $H{=}96$, $5\%$ at $H{=}336$ and $8\%$ at
$H{=}720$: in this matched comparison the value-space term accounts for the measured gain, and its
effect grows with horizon. Prior anchoring \citep{l2sp} is inert at the 32-frame clip length of the
main model, though it matters at 16 frames (Appendix~\ref{app:16f}).

Under this objective the budget curve is monotonically improving: $0.250 \to 0.245 \to 0.242$ at
$H{=}96$ for 5k, 20k and 60k updates. Scaling the other two knobs did not help --- a larger effective
batch is equal or worse at every horizon and a $4\times$ larger in-domain corpus ties ours ---
Appendix~\ref{app:failed}.

\subsection{Why naive continual pretraining fails}
\label{sec:diagnosis}

% Figure 1: the budget curve under the two objectives.
% Numbers: PAPER_MATERIAL.md sections 7-8 (VIDEO_TS_RESCUE_RESULTS.md 1.10-1.13).
\begin{figure}[t]
\centering
\begin{tikzpicture}

\begin{axis}[
  name=left,
  width=5.7cm, height=3.7cm,
  xmode=log, log basis x=10,
  xlabel={continual-pretraining steps},
  ylabel={held-out pixel loss},
  ylabel style={color=pixblue},
  yticklabel style={color=pixblue},
  ymin=0.0295, ymax=0.0382,
  xmin=3500, xmax=260000,
  xtick={5000,20000,60000,176000},
  xticklabels={5k,20k,60k,176k},
  ytick={0.031,0.033,0.035,0.037},
  yticklabel style={/pgf/number format/fixed, /pgf/number format/precision=4},
  scaled y ticks=false,
  title={\small (a) native masked-pixel objective},
  axis y line*=left, axis x line*=bottom,
  tick label style={font=\scriptsize}, label style={font=\scriptsize},
  every axis plot/.append style={thick},
]
\addplot[color=pixblue, mark=*, mark size=1.6pt]
  coordinates {(5000,0.0366) (20000,0.0333) (60000,0.0329) (176000,0.0312)};
\node[font=\scriptsize, color=pixblue, anchor=west] at (axis cs:23000,0.0361) {objective improves};
\end{axis}

\begin{axis}[
  width=5.7cm, height=3.7cm,
  at={(left.south east)}, anchor=south east,
  xmode=log, log basis x=10,
  ylabel={forecast MSE}, ylabel style={yshift=-2pt},
  ylabel style={color=fcred},
  yticklabel style={color=fcred},
  ymin=0.330, ymax=0.390,
  xmin=3500, xmax=260000,
  axis y line*=right, axis x line=none,
  ytick={0.34,0.35,0.36,0.37,0.38},
  tick label style={font=\scriptsize}, label style={font=\scriptsize},
]
\addplot[color=fcred, mark=square*, mark size=1.5pt, dashed, thick]
  coordinates {(5000,0.339) (20000,0.344) (176000,0.381)};
\node[font=\scriptsize, color=fcred, anchor=north east] at (axis cs:190000,0.372) {forecasting degrades};
\end{axis}

\begin{axis}[
  width=5.7cm, height=3.7cm,
  at={(left.south east)}, anchor=south west, xshift=2.35cm,
  xmode=log, log basis x=10,
  xlabel={continual-pretraining steps},
  ylabel={5-dataset mean MSE},
  xmin=3500, xmax=90000, ymin=0.230, ymax=0.395,
  xtick={5000,20000,60000}, xticklabels={5k,20k,60k},
  title={\small (b) forecast-aligned value-space objective},
  axis lines*=left,
  tick label style={font=\scriptsize}, label style={font=\scriptsize},
  legend style={font=\scriptsize, draw=none, fill=none, at={(0.99,0.30)}, anchor=east},
  every axis plot/.append style={thick, mark size=1.6pt},
]
\addplot[color=fcred, mark=triangle*] coordinates {(5000,0.384) (20000,0.379) (60000,0.377)};
\addlegendentry{$H{=}720$}
\addplot[color=black!70, mark=square*] coordinates {(5000,0.337) (20000,0.330) (60000,0.329)};
\addlegendentry{$H{=}336$}
\addplot[color=pixblue, mark=*] coordinates {(5000,0.250) (20000,0.245) (60000,0.242)};
\addlegendentry{$H{=}96$}
\end{axis}

\end{tikzpicture}
\caption{The same backbone and corpus under two objectives. \textbf{(a)} VideoMAE's native
masked-pixel objective: its own held-out loss falls monotonically while forecast error rises
monotonically over a $35\times$ range of budget. The 60k checkpoint was not evaluated downstream.
\textbf{(b)} The forecast-aligned value-space objective of Section~\ref{sec:objective}: the curve is
monotonically improving at every horizon. Panel (a) uses 16-frame clips and panel (b) 32-frame
clips; the ablation in Table~\ref{tab:ablation} compares the two objectives at matched clip length.}
\label{fig:budget}
\end{figure}


Figure~\ref{fig:budget}a is the observation that motivated the design. Continuing VideoMAE's native
masked-pixel objective on the rendered corpus improves its held-out pixel loss monotonically across
the four checkpoints we evaluated ($0.0366 \to 0.0312$) while forecast error rises monotonically
across the three we evaluated downstream ($0.339 \to 0.381$). Per cell, the 176k-step model is
$5$--$60\%$ worse than the 20k-step model, with the damage concentrated at long horizons; ETTh2 at
$H{=}720$ moves from $0.429$ to $0.617$. We did not evaluate the 60k native-objective checkpoint
downstream, so that point is missing rather than interpolated, and we describe this as a trend over
the checkpoints we measured rather than as a scaling law.

Alternative explanations were excluded by re-evaluation under matched software, and the training
loss is flat over the range, so this is not pixel-space overfitting; what does move monotonically is
the weight norm, which shrinks by roughly $30\%$ under decoupled weight decay over the $9\times$
longer schedule. Appendix~\ref{app:diagnosis} gives the checks.

Figure~\ref{fig:budget}b shows the corresponding curve under our objective. The panels differ in clip
length as well as objective (16 frames against 32), so they are not a controlled comparison; the
controlled one is Table~\ref{tab:ablation}. What the figure establishes is the direction of each
curve.

\subsection{Robustness and inference cost}
\label{sec:seeds}

Two seeds of the 20k model agree to within $0.002$ on the five-dataset means and on five of six
datasets of the full grid, ETTm1 being the only visible spread ($0.351$ against $0.360$); no second
seed of the 60k model exists, so the headline number carries no error bar. Inference is not
single-pass everywhere, since the rule of Section~\ref{sec:inference} averages frame scales on
short-horizon cells and rolls out beyond one prediction block. Appendix~\ref{app:inference} gives
the per-cell pass counts, and Appendix~\ref{app:failed} the variants that did not improve on the
reported system, including a two-seed prediction ensemble.
